\chapter{Phân tích và thiết kế}

\section{Kiến trúc tổng thể}

Hệ thống được tổ chức theo mô hình monorepo tách biệt client và server. Phía frontend gồm ba ứng dụng React độc lập cho ba vai trò. Phía backend dùng Express và được chia module theo nghiệp vụ. Cách tổ chức này giúp nhóm có thể phát triển theo nhánh vai trò, giảm xung đột khi nhiều thành viên cùng làm việc.

\begin{figure}[H]
\centering
\diagramimage{images/diagrams/architecture-overview.drawio.png}
\caption{Kiến trúc tổng quan hệ thống}
\end{figure}

\section{Cấu trúc frontend}

\begin{itemize}[leftmargin=1.2cm]
    \item \texttt{application/client/customer}: giao diện khách hàng, gồm trang chủ, tìm kiếm/lọc voucher, xem chi tiết, giỏ hàng, checkout, trạng thái thanh toán, hồ sơ cá nhân, ví E-Voucher/QR, lịch sử đơn hàng, đánh giá voucher, gửi khiếu nại, đăng ký khách hàng/đối tác, trang hỗ trợ/chính sách và chatbot.
    \item \texttt{application/client/partner}: cổng đối tác, gồm đăng nhập/đăng ký đối tác, dashboard kinh doanh, quản lý hồ sơ và chi nhánh, tạo/cập nhật/gửi duyệt/tạm ngưng voucher, báo cáo, tra cứu và xác thực voucher code bằng mã hoặc QR, chatbot và cập nhật realtime.
    \item \texttt{application/client/admin}: cổng quản trị, gồm dashboard, duyệt đối tác, quản lý người dùng, quản lý voucher, quản lý đơn hàng, xử lý khiếu nại, quản lý nội dung, nhật ký hệ thống và các chỉ số thống kê.
\end{itemize}

\section{Cấu trúc backend}

Backend được chia thành các module:

\begin{table}[H]
\centering
\small
\begin{tabularx}{\textwidth}{|p{3.6cm}|Y|}
\hline
\textbf{Module} & \textbf{Chức năng} \\ \hline
\texttt{auth} & Đăng ký, đăng nhập, JWT, hồ sơ, đổi mật khẩu, quên mật khẩu, OTP/email mô phỏng. \\ \hline
\texttt{shared/voucher} & API voucher công khai: danh sách voucher đã duyệt, danh mục, đối tác, tìm kiếm/lọc và xem chi tiết voucher. \\ \hline
\texttt{shared/chatbot} & API chatbot hỗ trợ người dùng trên các giao diện customer/partner. \\ \hline
\texttt{shared/event} & Kênh realtime bằng Server-Sent Events cho giao diện Quản trị viên và Đối tác để cập nhật trạng thái voucher, đơn hàng hoặc mã voucher. \\ \hline
\texttt{customer/order} & Kiểm tra giỏ hàng, checkout, cổng thanh toán mô phỏng/sandbox, trạng thái đơn, phát hành E-Voucher, lịch sử đơn, ví E-Voucher và đánh giá sau mua. \\ \hline
\texttt{customer/complaint} & Gửi khiếu nại, xem danh sách và chi tiết khiếu nại của khách hàng. \\ \hline
\texttt{partner} & Dashboard đối tác, quản lý chi nhánh, CRUD voucher, gửi duyệt/tạm ngưng voucher, báo cáo, kiểm tra và redeem voucher code theo phạm vi đối tác. \\ \hline
\texttt{admin} & Duyệt/từ chối đối tác, quản lý người dùng, quản lý đơn hàng tổng quát, xử lý khiếu nại, quản lý nội dung, nhật ký hệ thống và dữ liệu dashboard. \\ \hline
\texttt{admin/Voucher} & Danh sách voucher theo trạng thái, duyệt, từ chối, tạm ẩn/hiện voucher và thống kê voucher theo đối tác. \\ \hline
\texttt{admin/Order} & Danh sách đơn hàng, chi tiết đơn, xác nhận thanh toán và hoàn tiền mô phỏng. \\ \hline
\end{tabularx}
\caption{Module backend}
\end{table}

\section{Nguyên tắc thiết kế module}

Mỗi module backend tuân theo mô hình ba lớp:

\begin{enumerate}[leftmargin=1.2cm]
    \item Route: định nghĩa endpoint và middleware.
    \item Controller: nhận request, gọi service và trả response.
    \item Service: chứa logic nghiệp vụ và truy vấn database.
\end{enumerate}

Ví dụ, nhóm route quản trị voucher được mount ở \texttt{/api/admin/vouchers}; nhóm route đối tác được mount ở \texttt{/api/partner}; nhóm API chung được mount ở \texttt{/api/vouchers}; nhóm mua hàng được mount ở \texttt{/api/orders}; nhóm khiếu nại được mount ở \texttt{/api/complaints}; nhóm realtime được mount ở \texttt{/api/events}.

\section{Các sơ đồ đặc tả phần mềm}

Các sơ đồ dưới đây được dùng để đặc tả hệ thống ở nhiều góc nhìn: phạm vi chức năng, luồng nghiệp vụ, tương tác giữa thành phần, triển khai và mô hình dữ liệu nghiệp vụ. Nhóm ưu tiên các sơ đồ phản ánh đúng hiện trạng mã nguồn hiện tại thay vì mô tả các chức năng chưa có.

\subsection{System context diagram}

\subsubsection{Ranh giới và tác nhân ngoài hệ thống}

\begin{figure}[H]
\centering
\diagramimage{images/diagrams/system-context.drawio.png}
\caption{System context diagram xác định ranh giới và tác nhân ngoài hệ thống}
\end{figure}

\subsection{Use case diagram}

\subsubsection{Use case tổng quan hệ thống Dealzy}

\begin{figure}[H]
\centering
\diagramimage{images/diagrams/use-case-overview.drawio.png}
\caption{Use case tổng quan hệ thống Dealzy}
\end{figure}

\subsection{Activity diagram}

\subsubsection{Luồng mua và phát hành E-Voucher}

\begin{figure}[H]
\centering
\diagramimage{images/diagrams/activity-purchase.drawio.png}
\caption{Activity diagram cho quá trình mua và phát hành E-Voucher}
\end{figure}

Sơ đồ mô tả luồng từ khi Khách hàng chọn voucher, hệ thống kiểm tra điều kiện và tồn kho, tạo đơn hàng, tiếp nhận kết quả thanh toán đến khi phát hành E-Voucher. Trong phiên bản hiện tại, luồng này được hiện thực bằng API Order, transaction PostgreSQL và các cổng thanh toán mô phỏng/sandbox để phục vụ demo học thuật.

\subsubsection{Luồng tiếp nhận và xử lý khiếu nại}

\begin{figure}[H]
\centering
\diagramimage{images/diagrams/activity-complaint-processing.drawio.png}
\caption{Activity diagram cho tiếp nhận, xử lý và theo dõi khiếu nại}
\end{figure}

\subsection{Sequence diagram}

\subsubsection{Luồng mua và sử dụng voucher}

\begin{figure}[H]
\centering
\diagramimage{images/diagrams/sequence-purchase.drawio.png}
\caption{Sequence diagram cho mua, phát hành và xác thực E-Voucher}
\end{figure}

Trình tự tương tác thể hiện vai trò của giao diện Khách hàng, Order API, cổng thanh toán, PostgreSQL và giao diện Đối tác. Order API là thành phần điều phối việc kiểm tra dữ liệu, ghi nhận đơn hàng, cập nhật trạng thái và sinh E-Voucher sau khi thanh toán thành công để Đối tác tra cứu và xác nhận sử dụng.

\subsubsection{Xác thực và khôi phục tài khoản}

\begin{figure}[H]
\centering
\diagramimage{images/diagrams/sequence-auth-recovery.drawio.png}
\caption{Sequence diagram cho đăng ký, đăng nhập, JWT và khôi phục mật khẩu}
\end{figure}

\subsubsection{Tạo và duyệt voucher}

\begin{figure}[H]
\centering
\diagramimage{images/diagrams/sequence-voucher-approval.drawio.png}
\caption{Sequence diagram cho vòng đời tạo, gửi duyệt và phê duyệt voucher}
\end{figure}

\subsubsection{Kiểm tra và sử dụng E-Voucher}

\begin{figure}[H]
\centering
\diagramimage{images/diagrams/sequence-evoucher-redem.drawio.png}
\caption{Sequence diagram cho tra cứu và xác nhận sử dụng E-Voucher}
\end{figure}

\subsubsection{Cập nhật realtime}

\begin{figure}[H]
\centering
\diagramimage{images/diagrams/sequence-realtime-sse.drawio.png}
\caption{Sequence diagram cho kênh cập nhật realtime bằng Server-Sent Events}
\end{figure}

\subsection{State diagram}

\subsubsection{Vòng đời voucher và E-Voucher}

\begin{figure}[H]
\centering
\diagramimage{images/diagrams/state-voucher-lifcycle.drawio.png}
\caption{State diagram vòng đời voucher và E-Voucher}
\end{figure}

\subsubsection{Vòng đời đơn hàng}

\begin{figure}[H]
\centering
\diagramimage{images/diagrams/state-order-lifecyle.drawio.png}
\caption{State diagram vòng đời đơn hàng và các nhánh hoàn tồn kho}
\end{figure}

Đơn hàng bắt đầu ở trạng thái \texttt{Pending}. Khi nhận được xác nhận thanh toán hợp lệ, đơn chuyển sang \texttt{Paid}; nếu bị hủy, thanh toán thất bại hoặc quá hạn, đơn lần lượt chuyển sang \texttt{Cancelled}, \texttt{Failed} hoặc \texttt{Expired}. Trạng thái \texttt{Refunded} biểu diễn trường hợp hoàn tiền sau thanh toán. Các nhánh thất bại cần đi kèm cơ chế hoàn tồn kho để duy trì tính nhất quán dữ liệu.

\subsubsection{Vòng đời khiếu nại}

\begin{figure}[H]
\centering
\diagramimage{images/diagrams/state-complaint-lifecyle.drawio.png}
\caption{State diagram vòng đời khiếu nại}
\end{figure}

\subsection{Component diagram}

\subsubsection{Các thành phần của hệ thống}

\begin{figure}[H]
\centering
\diagramimage{images/diagrams/component-diagram.drawio.png}
\caption{Component diagram của hệ thống}
\end{figure}

\subsection{Deployment diagram}

\subsubsection{Mô hình triển khai demo}

\begin{figure}[H]
\centering
\diagramimage{images/diagrams/deployment-diagram.drawio.png}
\caption{Deployment diagram mức triển khai demo}
\end{figure}

\subsection{Ma trận bao phủ sơ đồ và phân hệ}

Để bảo đảm bộ sơ đồ phản ánh đầy đủ hệ thống hiện tại, bảng dưới đây ánh xạ từng phân hệ đang có trong mã nguồn đến sơ đồ đặc tả tương ứng. Một phân hệ chỉ được xem là đã bao phủ khi thể hiện được ít nhất một trong ba góc nhìn: chức năng, tương tác hoặc trạng thái dữ liệu.

\begin{longtable}{|p{3.2cm}|p{4.0cm}|p{6.2cm}|}
\hline
\textbf{Phân hệ hiện tại} & \textbf{Thành phần hiện thực} & \textbf{Sơ đồ thể hiện} \\ \hline
\endfirsthead
\hline
\textbf{Phân hệ hiện tại} & \textbf{Thành phần hiện thực} & \textbf{Sơ đồ thể hiện} \\ \hline
\endhead
Xác thực và hồ sơ & \texttt{auth}, đăng ký Khách hàng/Đối tác, hồ sơ & Use case tổng quan; sequence xác thực và khôi phục tài khoản \\ \hline
Khám phá voucher & \texttt{shared/voucher}, tìm kiếm và xem chi tiết voucher & Use case tổng quan; component diagram \\ \hline
Mua hàng và thanh toán & \texttt{customer/order}, MoMo, PayPal, VietQR & Activity mua E-Voucher; sequence mua/sử dụng; state diagram đơn hàng \\ \hline
Ví E-Voucher và đánh giá & \texttt{customer/order}, \texttt{E\_Vouchers}, \texttt{Reviews} & Sequence mua/sử dụng; state voucher \\ \hline
Khiếu nại & \texttt{customer/complaint}, API xử lý khiếu nại của Quản trị viên & Activity xử lý khiếu nại; state diagram khiếu nại \\ \hline
Vận hành đối tác & \texttt{partner}, dashboard, report, branch, redeem & Use case tổng quan; sequence tạo/duyệt voucher; sequence redeem \\ \hline
Quản trị hệ thống & \texttt{admin}, \texttt{admin/Voucher}, \texttt{admin/Order} & Use case tổng quan; sequence duyệt voucher; state đơn hàng \\ \hline
Realtime, chatbot và tiện ích ngoài & \texttt{shared/event}, \texttt{shared/chatbot}, email/SMS & Component; deployment; sequence realtime; context diagram \\ \hline
Nội dung và nhật ký & \texttt{Content\_Items}, \texttt{System\_Logs} & Use case tổng quan \\ \hline
\caption[]{Ma trận truy vết giữa phân hệ hiện tại và sơ đồ đặc tả}\\
\end{longtable}

\section{Luồng xử lý chính}

\subsection{Tạo và duyệt voucher}

Đối tác tạo voucher ở trạng thái Pending. Quản trị viên xem danh sách voucher Pending, kiểm tra thông tin, sau đó duyệt sang Approved, từ chối sang Rejected kèm lý do hoặc chuyển sang Suspended để tạm ẩn.

\subsection{Tìm kiếm và xem voucher}

Khách hàng chỉ nhìn thấy voucher có trạng thái Approved và còn hiệu lực. API tìm kiếm hỗ trợ lọc theo từ khóa, danh mục, giá, mức giảm, khu vực và đối tác.

\subsection{Xác thực voucher code}

Đối tác nhập voucher code. Backend kiểm tra code có tồn tại, thuộc voucher của đối tác, trạng thái còn Unused, chưa hết hạn và chi nhánh xác thực thuộc đúng đối tác. Nếu hợp lệ, hệ thống cập nhật trạng thái E\_Voucher thành Used và ghi nhận thời điểm sử dụng.

\section{Thiết kế phân quyền}

Backend sử dụng JWT. Middleware \texttt{authMiddleware} đọc token từ header Authorization. Các route admin và partner bổ sung kiểm tra role trước khi cho phép thao tác. Thiết kế này giúp tách rõ quyền truy cập:

\begin{itemize}[leftmargin=1.2cm]
    \item Quản trị viên chỉ truy cập endpoint quản trị.
    \item Đối tác chỉ truy cập endpoint đối tác.
    \item API voucher công khai chỉ trả voucher đã được duyệt.
\end{itemize}

\section{Thiết kế cơ sở dữ liệu}

\subsection{Tổng quan}

Cơ sở dữ liệu sử dụng PostgreSQL và được định nghĩa trong file \texttt{scriptDatabase.sql}. Script bao gồm lệnh xóa bảng cũ, tạo bảng, tạo trigger/ràng buộc nghiệp vụ và chèn dữ liệu mẫu.

\subsection{Các bảng chính}

\begin{longtable}{|p{3.4cm}|p{9.8cm}|}
\hline
\textbf{Bảng} & \textbf{Vai trò} \\ \hline
\endfirsthead
\hline
\textbf{Bảng} & \textbf{Vai trò} \\ \hline
\endhead
Users & Thông tin đăng nhập, email, phone, role, token khôi phục mật khẩu. \\ \hline
Customers & Hồ sơ khách hàng, ngày sinh, địa chỉ, trạng thái hoạt động. \\ \hline
Partners & Hồ sơ doanh nghiệp, người đại diện, mã số thuế, trụ sở, trạng thái duyệt. \\ \hline
Branches & Danh sách chi nhánh của đối tác. \\ \hline
Categories & Danh mục voucher. \\ \hline
Vouchers & Thông tin voucher, giá, số lượng, trạng thái, thời gian, điều kiện, chính sách. \\ \hline
Voucher\_Branches & Quan hệ nhiều-nhiều giữa voucher và chi nhánh áp dụng. \\ \hline
Orders & Đơn hàng mẫu, người mua, tổng tiền, trạng thái. \\ \hline
Order\_Items & Chi tiết đơn hàng, voucher, số lượng, giá tại thời điểm mua. \\ \hline
E\_Vouchers & Mã voucher điện tử, trạng thái sử dụng, ngày phát hành, ngày hết hạn, chi nhánh sử dụng. \\ \hline
Reviews & Đánh giá voucher của khách hàng. \\ \hline
Complaints & Khiếu nại/phản hồi của người dùng. \\ \hline
System\_Logs & Bảng dự kiến lưu vết hành động quan trọng. \\ \hline
\caption[]{Các bảng dữ liệu chính}\\
\end{longtable}

\subsection{ERD}

\begin{figure}[H]
\centering
\diagramimage{images/diagrams/database-erd.drawio.png}
\caption{ERD hệ thống voucher}
\end{figure}

ERD được tổ chức thành các nhóm dữ liệu chính:
\begin{itemize}[leftmargin=1.2cm]
    \item Nhóm người dùng gồm Users, Customers và Partners, lưu tài khoản và hồ sơ theo từng vai trò.
    \item Nhóm voucher gồm Vouchers, Categories, Branches và Voucher\_Branches, mô tả sản phẩm voucher cùng phạm vi áp dụng.
    \item Nhóm giao dịch gồm Orders, Order\_Items và E\_Vouchers, hỗ trợ mô hình hóa quy trình mua và phát hành mã điện tử.
    \item Nhóm phản hồi gồm Reviews, Complaints, Complaint\_Vouchers và Complaint\_Responses, lưu đánh giá sau giao dịch và hỗ trợ xử lý khiếu nại.
\end{itemize}

\subsection{Ràng buộc nghiệp vụ trong database}

Một số quy tắc được hiện thực trực tiếp trong schema hoặc trigger:

\begin{itemize}[leftmargin=1.2cm]
    \item \texttt{chk\_price}: giá bán phải nhỏ hơn giá gốc.
    \item \texttt{chk\_stock}: tồn kho không âm và không vượt quá tổng phát hành.
    \item \texttt{chk\_dates}: ngày hết hạn phải sau ngày bắt đầu.
    \item \texttt{trg\_validate\_order\_item}: kiểm tra voucher đã duyệt, còn hạn và còn tồn khi thêm order item.
    \item \texttt{trg\_validate\_voucher\_usage}: kiểm tra chi nhánh xác thực thuộc đúng đối tác của voucher.
    \item \texttt{trg\_validate\_review}: chỉ cho phép đánh giá khi khách hàng có đơn liên quan đến voucher.
\end{itemize}

\subsection{Chỉ mục dữ liệu}

Qua kiểm tra metadata PostgreSQL trên Supabase, hệ thống hiện có nhiều chỉ mục thuộc các schema hạ tầng của Supabase như \texttt{auth}, \texttt{storage} và \texttt{realtime}. Các chỉ mục này phục vụ cơ chế xác thực, lưu trữ tệp và realtime của nền tảng Supabase, không phải thành phần thiết kế nghiệp vụ riêng của ứng dụng Dealzy. Trong phạm vi schema nghiệp vụ \texttt{public}, ngoài các chỉ mục được PostgreSQL tự động tạo từ khóa chính và ràng buộc UNIQUE, dự án có hai chỉ mục tùy chỉnh sau:

\begin{table}[H]
\centering
\footnotesize
\begin{tabularx}{\textwidth}{|p{3.2cm}|p{3.0cm}|Y|}
\hline
\textbf{Chỉ mục} & \textbf{Đối tượng áp dụng} & \textbf{Ý nghĩa và tác dụng} \\ \hline
\shortstack[l]{\texttt{idx\_users\_email}\\\texttt{\_unique\_not\_null}} & \texttt{users.email} & Bảo đảm email là duy nhất theo dạng không phân biệt chữ hoa/thường đối với các bản ghi có email. Chỉ mục này giúp kiểm soát trùng lặp khi đăng ký tài khoản, đồng thời hỗ trợ tra cứu nhanh trong các luồng đăng nhập, xác minh hoặc khôi phục tài khoản bằng email. \\ \hline
\shortstack[l]{\texttt{idx\_users\_phone}\\\texttt{\_unique\_not\_null}} & \texttt{users.phone} & Bảo đảm số điện thoại là duy nhất khi người dùng có khai báo số điện thoại, nhưng vẫn cho phép các tài khoản chưa cập nhật số điện thoại. Chỉ mục này hỗ trợ kiểm tra trùng số điện thoại và các luồng xác thực OTP/SMS mô phỏng trong hệ thống. \\ \hline
\end{tabularx}
\caption{Các chỉ mục custom trong schema nghiệp vụ}
\end{table}

Về mặt cài đặt, hai chỉ mục trên đều là partial unique index. Chỉ mục email áp dụng trên biểu thức \texttt{lower(email)} với điều kiện email khác \texttt{NULL}; chỉ mục số điện thoại áp dụng trên cột \texttt{phone} với điều kiện số điện thoại khác \texttt{NULL}. Cách thiết kế này phù hợp với dữ liệu người dùng vì email và số điện thoại có thể chưa được khai báo ở thời điểm tạo tài khoản, nhưng khi đã có giá trị thì cần bảo đảm tính duy nhất.

Ngoài hai chỉ mục tùy chỉnh trên, PostgreSQL tự tạo chỉ mục cho các khóa chính của những bảng nghiệp vụ như người dùng, voucher, đơn hàng và E-Voucher. Các ràng buộc UNIQUE trên username, mã E-Voucher, quan hệ đánh giá của khách hàng, tên danh mục và khóa nội dung cấu hình cũng được ánh xạ thành chỉ mục tương ứng. Nhóm chỉ mục tự sinh này vừa bảo đảm toàn vẹn dữ liệu, vừa tăng tốc tra cứu theo khóa định danh, mã voucher điện tử, username, danh mục và nội dung cấu hình.

\subsection{Dữ liệu mẫu}

Seed data hiện có:

\begin{itemize}[leftmargin=1.2cm]
    \item Tài khoản admin, đối tác và khách hàng với password hash.
    \item Sáu đối tác mẫu: Sheraton, Fantastic Travel, Glow Spa, Nike, Hokkaido Sushi, CGV.
    \item Danh mục Dining, Travel, Beauty, Shopping, Entertainment.
    \item Nhiều voucher mẫu có trạng thái, điều kiện áp dụng, chính sách hoàn/hủy và chi nhánh.
    \item Đơn hàng, order item, E\_Voucher và review mẫu để minh họa nghiệp vụ.
\end{itemize}

